% Part: lambda-calculus
% Chapter: lambda-definability
% Section: minimization

\documentclass[../../../include/open-logic-section]{subfiles}

\begin{document}

\olfileid[fr]{lam}{rep}{min}
\olsection{Minimisation}

Les fonctions récursives générales
s'obtiennent à partir des fonctions
initiales $\Zero$, $\Succ$ et
$\Proj{n}{i}$ par composition,
récurrence primitive et minimisation
régulière. Pour montrer qu'elles sont
toutes $\lambda$-définissables, il
reste à établir que la minimisation
régulière d'une fonction
$\lambda$-définissable est encore
$\lambda$-définissable.

\begin{lem}
\ollabel{lem:min}
Si $f(x_1,\dots,x_k,y)$ est régulière
et $\lambda$-définissable, alors la
fonction~$g$ donnée par
\[
  g(x_1,\dots,x_k)
    =\umin{y}{f(x_1,\dots,x_k,y)=0}
\]
est aussi $\lambda$-définissable.
\end{lem}

\begin{proof}
Supposons que le lambda-terme~$F$
définisse la fonction régulière
$f(\vec x,y)$. Pour définir~$g$,
utilisons une fonction de recherche
et un combinateur de point fixe.
Posons
\begin{align*}
  \fn{Search} &\ident
    \lambd[g][\lambd[f\,\vec x\,y][
      \fn{IsZero}(f\,\vec x\,y)\,y\,
      (g\,f\,\vec x\,(\fn{Succ}\,y))]],\\
  H &\ident
    \lambd[\vec x][
      (Y\,\fn{Search})F\,\vec x\,\num{0}],
\end{align*}
où~$Y$ est le combinateur de Turing
de la \olref[fp]{defn:Turing-Y}.
À partir de~$y$, le terme
$\fn{Search}$ teste si
$f(\vec x,y)=0$. Si oui, il renvoie~$y$;
sinon, il reprend la recherche à
$\fn{Succ}\,y$ en conservant le même
terme~$f$. Ainsi,
$(Y\,\fn{Search})F\,
\num{n_1}\dots\num{n_k}\,\num{0}$
renvoie le plus petit~$m$ pour lequel
$f(n_1,\dots,n_k,m)=0$.

Plus précisément, si
$f(n_1,\dots,n_k,m)=0$, alors
\[
  (Y\,\fn{Search})F\,
    \num{n_1}\dots\num{n_k}\,\num{m}
    \red\num{m}.
\]
Si la valeur de~$f$ est non nulle,
on obtient
\[
  (Y\,\fn{Search})F\,
    \num{n_1}\dots\num{n_k}\,\num{m}
    \red (Y\,\fn{Search})F\,
      \num{n_1}\dots\num{n_k}\,\num{m+1}.
\]
La régularité de~$f$ signifie ici
que~$f$ est totale et que, pour
chaque tuple $\vec n$, il existe un
$y$ tel que $f(\vec n,y)=0$.
Toutes les étapes précédant le
premier zéro sont donc définies et
la recherche l'atteint. Par conséquent,
\[
  (Y\,\fn{Search})F\,
    \num{n_1}\dots\num{n_k}\,\num{0}
    \red
    \num{g(n_1,\dots,n_k)}.
\]
Le terme~$H$ définit bien~$g$.
\end{proof}

\begin{prop}
Toute fonction récursive générale
est $\lambda$-définissable.
\end{prop}

\begin{proof}
Les fonctions initiales sont
$\lambda$-définissables par le
\olref[prf]{lem:basic}.
Le \olref[prf]{lem:comp},
le \olref[prf]{lem:prim} et le
\olref{lem:min} donnent la fermeture
par composition, récurrence primitive
et minimisation régulière.
\end{proof}

La recherche maintient les paramètres $\vec x$ et la fonction
$F$ à chaque appel. Seule la valeur de l'indice change.
La totalité garantit un résultat pour chaque test et
l'existence d'un zéro garantit que l'un d'eux arrête la recherche.

\end{document}
